\ifdefined\gkver\else\def\gkver{student}\fi
\documentclass[\gkver]{gaokaozhenti}
\begin{document}

\chapter{2010年广东理综}
%% paperType: 理综物理部分

\begin{enumerate}
\item
%% number: 13
%% paperName: 2010年广东理综第 13 题 4 分
%% typeId: 1
%% type: 单选题
%% chapter: 力物体的平衡
%% point: 三力平衡
%% method: 无
%% score: 4
%% degree: 650
%% duplicateId: 0
%% body:
如图为节日里悬挂灯笼的一种方式，A、B 点等高，O 为结点，轻绳 AO、BO 长度相等，拉力分别为 $F_{A}$、$F_{B}$，灯笼受到的重力为 $G$。下列表述正确的是\xzanswer{B}

\onepicture[w=4.2cm]{figs/13.png}

\fourchoices[answer=B]
{$F_{A}$ 一定小于 $G$}
{$F_{A}$ 与 $F_{B}$ 大小相等}
{$F_{A}$ 与 $F_{B}$ 是一对平衡力}
{$F_{A}$ 与 $F_{B}$ 大小之和等于 $G$}

%% answer: B
%% memo:
\memoanswer{
由等高等长知，左右力对称，选项 B 正确。选项 A 错误，有可能大于；选项 D 错误，不是大小之和而是矢量之和。选项 C 错误，这两个力的矢量和与重力是平衡力。
}

\item
%% number: 14
%% paperName: 2010年广东理综第 14 题 4 分
%% typeId: 1
%% type: 单选题
%% chapter: 气体性质
%% point: 热力学第一定律
%% method: 无
%% score: 4
%% degree: 650
%% duplicateId: 0
%% body:
如图是密闭的气缸，外力推动活塞 P 压缩气体，对缸内气体做功 $800\UJ$，同时气体向外界放热 $200\UJ$，缸内气体的\xzanswer{A}

\onepicture[w=2.8cm]{figs/14.png}

\fourchoices[answer=A]
{温度升高，内能增加 $600\UJ$}
{温度升高，内能减少 $200\UJ$}
{温度降低，内能增加 $600\UJ$}
{温度降低，内能减少 $200\UJ$}

%% answer: A
%% memo:
\memoanswer{
由能量守恒，$\Delta E=Q+W=-200+800=600\UJ$，内能增加 $600\UJ$，则温度一定升高。
}

\item
%% number: 15
%% paperName: 2010年广东理综第 15 题 4 分
%% typeId: 1
%% type: 单选题
%% chapter: 气体性质
%% point: 玻意耳定律
%% method: 无
%% score: 4
%% degree: 650
%% duplicateId: 0
%% body:
如图所示，某种自动洗衣机进水时，与洗衣缸相连的细管中会封闭一定质量的空气，通过压力传感器感知管中的空气压力，从而控制进水量。设温度不变，洗衣缸内水位升高，则细管中被封闭的空气\xzanswer{B}

\onepicture[w=5.4cm]{figs/15.png}

\fourchoices[answer=B]
{体积不变，压强变小}
{体积变小，压强变大}
{体积不变，压强变大}
{体积变小，压强变小}

%% answer: B
%% memo:
\memoanswer{
由图可知空气被封闭在细管内，水面升高时，根据玻意耳定律，气体压强增大，气体体积减小。
}

\item
%% number: 16
%% paperName: 2010年广东理综第 16 题 4 分
%% typeId: 1
%% type: 单选题
%% chapter: 电磁感应
%% point: 电磁感应电路分析
%% method: 无
%% score: 4
%% degree: 650
%% duplicateId: 0
%% body:
如图所示，平行导轨间有一矩形的匀强磁场区域，细金属棒 PQ 沿导轨从 MN 处匀速运动到 $M'N'$ 的过程中，棒上感应电动势 $E$ 随时间 $t$ 变化的图示，可能正确的是\xzanswer{A}

\onepicture[w=4.4cm]{figs/16a.png}

\fourchoices[answer=A, ispicture=true, h=1.6cm]
{figs/16b.png}
{figs/16c.png}
{figs/16d.png}
{figs/16e.png}

%% answer: A
%% memo:
\memoanswer{
导线做匀速直线运动切割磁感线时，$E=BLv$，是常数。开始没有切割，没有电动势，最后一段也没有切割，没有电动势。
}

\item
%% number: 17
%% paperName: 2010年广东理综第 17 题 6 分
%% typeId: 1
%% type: 单选题
%% chapter: 直线运动
%% point: 运动图象
%% method: 无
%% score: 6
%% degree: 450
%% duplicateId: 0
%% body:
如图所示是某质点运动的速度图像，由图像得到的正确结果是\xzanswer{BC}

\onepicture[w=5cm]{figs/17.png}

\fourchoices[answer=BC]
{$0\sim1\Us$ 内的平均速度是 $2\Ums$}
{$0\sim2\Us$ 内的位移大小是 $3\Um$}
{$0\sim1\Us$ 内的加速度大于 $2\sim4\Us$ 内的加速度}
{$0\sim1\Us$ 内的运动方向与 $2\sim4\Us$ 内的运动方向相反}

%% answer: BC
%% memo:
\memoanswer{
选项 A 错误，平均速度应该是 $1\Ums$。选项 D 错误，速度都为正，方向相同。
}

\item
%% number: 18
%% paperName: 2010年广东理综第 18 题 6 分
%% typeId: 2
%% type: 多选题
%% chapter: 原子和原子核
%% point: 人工核反应
%% method: 无
%% score: 6
%% degree: 650
%% duplicateId: 0
%% body:
关于核衰变和核反应的类型，下列表述正确的有\xzanswer{AC}

\fourchoices[answer=AC]
{$\ce{^{238}_{92}U -> ^{234}_{90}Th + ^{4}_{2}He}$ 是 $\alpha$ 衰变}
{$\ce{^{14}_{7}N + ^{4}_{2}He -> ^{17}_{8}O + ^{1}_{1}H}$ 是 $\beta$ 衰变}
{$\ce{^{2}_{1}H + ^{3}_{1}H -> ^{4}_{2}He + ^{1}_{0}n}$ 是轻核聚变}
{$\ce{^{82}_{34}Se -> ^{82}_{36}Kr + 2^{0}_{-1}e}$ 是重核裂变}

%% answer: AC
%% memo:
\memoanswer{
B 选项的核反应方程是卢瑟福发现质子的核反应方程，B 错误。选项 D 核反应方程是 $\beta$ 衰变，D 错误。
}

\item
%% number: 19
%% paperName: 2010年广东理综第 19 题 6 分
%% typeId: 2
%% type: 多选题
%% chapter: 交流电
%% point: 交流电
%% method: 无
%% score: 6
%% degree: 650
%% duplicateId: 0
%% body:
如图所示是某种正弦式交变电压的波形图，由图可确定该电压的\xzanswer{BC}

\onepicture[w=5.5cm]{figs/19.png}

\fourchoices[answer=BC]
{周期是 $0.01\Us$}
{最大值是 $311\UV$}
{有效值是 $220\UV$}
{表达式为 $U=220\sin100\pi t$（$\UV$）}

%% answer: BC
%% memo:
\memoanswer{
由题图可得周期为 $0.02\Us$，A 错误。选项 D 中峰值应该为 $311\UV$，即 $U=311\sin100\pi t$（$\UV$），D 错误。
}

\item
%% number: 20
%% paperName: 2010年广东理综第 20 题 6 分
%% typeId: 2
%% type: 多选题
%% chapter: 力物体的平衡
%% point: 力
%% method: 无
%% score: 6
%% degree: 650
%% duplicateId: 0
%% body:
下列关于力的说法正确的是\xzanswer{BD}

\fourchoices[answer=BD]
{作用力和反作用力作用在同一物体上}
{太阳系中的行星均受到太阳的引力作用}
{运行的人造地球卫星所受引力的方向不变}
{伽利略的理想实验说明了力不是维持物体运动的原因}

%% answer: BD
%% memo:
\memoanswer{
A 错误，应该在两个物体上。由于方向是变化的，C 错误。
}

\item
%% number: 21
%% paperName: 2010年广东理综第 21 题 6 分
%% typeId: 2
%% type: 多选题
%% chapter: 电场
%% point: 电场线
%% method: 无
%% score: 6
%% degree: 650
%% duplicateId: 0
%% body:
如图所示是某一点电荷的电场线分布图，下列表述正确的是\xzanswer{BD}

\onepicture[w=3.8cm]{figs/21.png}

\fourchoices[answer=BD]
{a 点的电势高于 b 点的电势}
{该点电荷带负电}
{a 点和 b 点电场强度的方向相同}
{a 点的电场强度大于 b 点的电场强度}

%% answer: BD
%% memo:
\memoanswer{
沿着电场线方向电势逐渐降低，a 点电势低于 b 点电势，A 错误。由题图可看出 a 点和 b 点电场强度的方向不同，C 错误。
}

\item
%% number: 34
%% paperName: 2010年广东理综第 34 题 6 分
%% typeId: 4
%% type: 实验题
%% chapter: 直线运动
%% point: 打点计时器公式
%% method: 无
%% score: 6
%% degree: 650
%% duplicateId: 0
%% body:
如图是某同学在做匀变速直线运动实验中获得的一条纸带。

\begin{enumerate}
	\item
	已知打点计时器电源频率为 $50\UHz$，则纸带上打相邻两点的时间间隔为\tkanswer[1.8cm]{$0.02\Us$}。

	\item
	ABCD 是纸带上四个计数点，每两个相邻计数点间有四个点没有画出。从图中读出 A、B 两点间距 $s=$\tkanswer[1.6cm]{$0.70\Ucm$}；C 点对应的速度是\tkanswer[2.2cm]{$0.100\Ums$}（计算结果保留三位有效数字）。
\end{enumerate}

\onepicture[w=12cm]{figs/34tape.png}

%% answer: 
\jdanswer{
\begin{enumerate}
	\item
	$0.02\Us$

	\item
	$0.70\Ucm$，$0.100\Ums$

\end{enumerate}
}
%% memo:
\memoanswer{
C 点的瞬时速度等于 BD 段的平均速度。
}

\item
%% number: 34
%% paperName: 2010年广东理综第 34 题 12 分
%% typeId: 4
%% type: 实验题
%% chapter: 电路
%% point: 测电动势与内阻
%% method: 无
%% score: 12
%% degree: 450
%% duplicateId: 0
%% body:
某同学利用电压表和电阻箱测定干电池的电动势和内阻，使用的器材还包括定值电阻（$R_{0}=5\UO$）一个，开关两个，导线若干，实验原理图如图（a）。

\threepicture[numA=（a）, numB=（b）, numC=（c）, labelprefix={}, wA=3.2cm, wB=6.2cm, wC=2.6cm]
{figs/34a.png}
{figs/34b.png}
{figs/34c.png}

\begin{enumerate}
	\item
	在图（b）的实物图中，已正确连接了部分电路，请完成余下电路的连接。

	\item
	请完成下列主要实验步骤；

	A、检查并调节电压表指针指零；调节电阻箱，示数如图（c）所示，读得电阻值是\tkanswer[1.6cm]{$20\UO$}；

	B、将开关 $S_{1}$ 闭合，开关 $S_{2}$ 断开，电压表的示数是 $1.49\UV$；

	C、将开关 $S_{2}$\tkanswer[1.2cm]{闭合}，电压表的示数是 $1.16\UV$；断开开关 $S_{1}$。

	\item
	使用测得的数据，计算出干电池的内阻是\tkanswer[1.6cm]{$0.69\UO$}（计算结果保留二位有效数字）。

	\item
	由于所有电压表不是理想电压表，所以测得的电动势比实际值偏\tkanswer[1.2cm]{小}（填“大”或“小”）。
\end{enumerate}

%% answer: 
\jdanswer{
\begin{enumerate}
	\item
	如图所示。

	\onepicture[w=8.5cm]{figs/34_an.png}

	\item
	$20\UO$，闭合

	\item
	$0.69\UO$

	\item
	小

\end{enumerate}
}
%% memo:
\memoanswer{
串联电路中，部分电路的电压与电阻成正比，$\frac{1.49-1.16}{5+r}=\frac{1.16}{20}$，可得 $r=0.69\UO$。
}

\item
%% number: 35
%% paperName: 2010年广东理综第 35 题 18 分
%% typeId: 6
%% type: 计算题
%% chapter: 运动和力
%% point: 动量和能量
%% method: 无
%% score: 18
%% degree: 450
%% duplicateId: 0
%% body:
如图所示，一条轨道固定在竖直平面内，粗糙的 ab 段水平，bcde 段光滑，cde 段是以 O 为圆心、$R$ 为半径的一小段圆弧。可视为质点的物块 A 和 B 紧靠在一起，静止于 b 处，A 的质量是 B 的 3 倍。两物块在足够大的内力作用下突然分离，分别向左、右始终沿轨道运动。B 到 d 点时速度沿水平方向，此时轨道对 B 的支持力大小等于 B 所受重力的 $\frac{3}{4}$，A 与 ab 段的动摩擦因数为 $\mu$，重力加速度 $g$，求：

\begin{enumerate}
	\item
	物块 B 在 d 点的速度大小；

	\item
	物块 A 滑行的距离 $s$。
\end{enumerate}

\onepicture[w=9cm]{figs/35.png}

%% answer: 
\jdanswer{
\begin{enumerate}
	\item
	$v=\frac{\sqrt{Rg}}{2}$

	\item
	$s=\frac{R}{8\mu}$

\end{enumerate}
}
%% memo:
\memoanswer{
（1）B 在 d 点，根据牛顿第二定律有：

$mg-\frac{3}{4}mg=m\frac{v^{2}}{R}$，解得：$v=\frac{\sqrt{Rg}}{2}$。

（2）B 从 b 到 d 过程，只有重力做功，机械能守恒有：

$\frac{1}{2}mv_{B}^{2}=mgR+\frac{1}{2}mv^{2}$……①

AB 分离过程动量守恒有：$3mv_{A}=mv_{B}$……②

A 匀减速直线运动，用动能定理得，$0-\frac{1}{2}\cdot3mv_{A}^{2}=-\mu\cdot3mgs$……③

联立①②③，解得：$s=\frac{R}{8\mu}$。
}

\item
%% number: 36
%% paperName: 2010年广东理综第 36 题 18 分
%% typeId: 6
%% type: 计算题
%% chapter: 磁场
%% point: 带电粒子在磁场中的运动
%% method: 无
%% score: 18
%% degree: 250
%% duplicateId: 0
%% body:
如图（a）所示，左为某同学设想的粒子速度选择装置，由水平转轴及两个薄盘 $N_{1}$、$N_{2}$ 构成，两盘面平行且与转轴垂直，相距为 $L$，盘上各开一狭缝，两狭缝夹角 $\theta$ 可调（如图（b））；右为水平放置的长为 $d$ 的感光板，板的正上方有一匀强磁场，方向垂直纸面向外，磁感应强度为 $B$。一小束速度不同、带正电的粒子沿水平方向射入 $N_{1}$，能通过 $N_{2}$ 的粒子经 O 点垂直进入磁场。O 到感光板的距离为 $d/2$，粒子电荷量为 $q$，质量为 $m$，不计重力。

\threepicture[numA=（a）, numB=（b）, numC=（c）, labelprefix={}, wA=6cm, wB=4.2cm, wC=3.6cm]
{figs/36a.png}
{figs/36b.png}
{figs/36c.png}

\begin{enumerate}
	\item
	若两狭缝平行且盘静止（如图（c）），某一粒子进入磁场后，竖直向下打在感光板中心点 M 上，求该粒子在磁场中运动的时间 $t$；

	\item
	若两狭缝夹角为 $\theta_{0}$，盘匀速转动，转动方向如图（b）。要使穿过 $N_{1}$、$N_{2}$ 的粒子均打到感光板 $P_{1}P_{2}$ 连线上。试分析盘转动角速度 $\omega$ 的取值范围（设通过 $N_{1}$ 的所有粒子在盘转一圈的时间内都能到达 $N_{2}$）。
\end{enumerate}

%% answer: 
\jdanswer{
\begin{enumerate}
	\item
	$t=\frac{\pi m}{2qB}$

	\item
	$\frac{qBd\theta_{0}}{4mL}\leq\omega\leq\frac{5qBd\theta_{0}}{4mL}$

\end{enumerate}
}
%% memo:
\memoanswer{
（1）粒子在磁场中匀速圆周运动，洛伦兹力提供向心力，$qvB=m\frac{v^{2}}{R}$，周长=周期×速度，$2\pi R=vT$，又 $t=\frac{T}{4}$，解得：$t=\frac{\pi m}{2qB}$。

（2）速度最小时，$L=v_{1}t_{1}$，$\theta_{0}=\omega_{1}t_{1}$，$qv_{1}B=m\frac{v_{1}^{2}}{\frac{d}{4}}$，解得：$\omega_{1}=\frac{qBd\theta_{0}}{4mL}$。

速度最大时，$L=v_{2}t_{2}$，$\theta_{0}=\omega_{2}t_{2}$，$qv_{2}B=m\frac{v_{2}^{2}}{R}$，$R^{2}=(R-\frac{d}{2})^{2}+d^{2}$，解得：$\omega_{2}=\frac{5qBd\theta_{0}}{4mL}$，所以 $\frac{qBd\theta_{0}}{4mL}\leq\omega\leq\frac{5qBd\theta_{0}}{4mL}$。
}

\end{enumerate}

\end{document}
